\section{Soul\'e's 2003 bound and the explicit constants}
\label{sec:soule}

This section records a source comparison; Soul\'e's theorem
is cited and is not reproved. Let $F$ have degree
$r=r_1+2r_2$ and absolute discriminant $D$, and let
$A=\mathcal O_F$. For a finite abelian group $G$,
$\operatorname{card}_b(G)$ denotes the largest divisor of
$|G|$ whose prime factors are all strictly greater than $b$.
Thus it removes the entire primary parts at primes at most $b$.

Soul\'e, \emph{A bound for the torsion in the K-theory of
algebraic integers}, Documenta Mathematica, Extra Volume Kato
(2003), 761--788, Theorem~2, \S4.2, printed p.~779, states
\[
 \operatorname{card}_{\max(r+1,2m+2)}
       K_m(\mathcal O_F)_{\mathrm{tors}}\leq k(m).
 \tag{SO.1}
\]
The definitions in \S\S3.2, 4.1--4.2, pp.~777--778, are
\[
\begin{aligned}
 N&=2m+1,&
 \widetilde w&=r_1N(N+1)/2+r_2N^2-N,\\
 k(m)&=h(m,N),&
 h(q,N)&=a(N)^{(N+1)c(\widetilde w-q-1,N)},\\
 c(j,N)&=\binom{a(N)}{N+j},&
 a(N)&=2^{N(r+3)}
                  \left(\prod_{i=1}^N T_i\right)^{r/2}.
\end{aligned}
 \tag{SO.2}
\]
The constants $T_i$ are those in Proposition~2,
pp.~774--775; this is the exact recursive constant used by
the theorem. Its fully explicit simplified upper bounds
are Proposition~4, \S4.4, pp.~779--780:
\[
\begin{array}{ll}
\text{general number fields }(r\geq2):&
 \log\log k(m)\leq220\,m^4\log m\,
                   r^{4r}\sqrt D(\log D)^{r-1},\\[2pt]
\text{class number one }(r\geq2):&
 \log\log k(m)\leq210\,m^4\log m\,
                   r^{4r}\sqrt D(\log D)^{r-1},\\[2pt]
\text{imaginary quadratic}:&
 \log\log k(m)\leq1120\,m^4\log m\log D,\\[2pt]
\text{imaginary quadratic, class number one}:&
 \log\log k(m)\leq510\,m^4\log m\log D.
\end{array}
 \tag{SO.3}
\]
The general estimates are used for $m\geq2$ as in their proof,
pp.~781--784. These coefficients are the integers
$220,210,1120,510$; they are not powers of two, eleven, or five.
For $F=\mathbf Q$, Proposition~4(iv), printed p.~780, gives
\[
\begin{aligned}
 \log\log k(m)&\leq8m^4\log m\qquad(m\geq9),\\
 \log\log k(7)&\leq40545,\qquad
 \log\log k(8)\leq70130.
\end{aligned}
 \tag{SO.4}
\]
The pages containing (SO.1), (SO.3), and (SO.4) were checked
visually against the attached PDF.

\subsection{Comparison in the degrees of interest}

For $m=4k$, $k\geq3$, Soul\'e gives
\[
 \log\left(
   \operatorname{card}_{8k+2} K_{4k}(\mathbf Z)
 \right)
 \leq \exp\!\bigl(8(4k)^4\log(4k)\bigr).
 \tag{SO.5}
\]
The inherited VK-EXP7 result instead bounds the complete order
by
\[
       \log|K_{4k}(\mathbf Z)|\leq C_7(2k)^7.
 \tag{SO.6}
\]
There are two separate improvements: polynomial growth for
the logarithm in (SO.6), and inclusion of every primary part.
Soul\'e's excluded-prime cutoff grows with the degree.
The verified Vandiver range removes that cutoff only in
degrees where all its primes fall in the verified range;
it must not be silently removed for every $k$.
The fresh KT-C7-EXPLICIT derivation below makes one admissible
archive-based constant explicit. KT-QUARTIC-BOUND gives the
sharper, independently derived bound
$\log|K_{4k}(\mathbf Z)|<\frac25(2k)^4$.

\paragraph{Arithmetic-geometric meaning.}
Soul\'e's Theorem~3, \S4.3, printed p.~779, also states
\[
 \prod_{\substack{p\geq4n-1\\p\geq r+2}}
 \left|H^2_{\acute et}(\operatorname{Spec}\mathcal O_F[1/p],
                                      \mathbf Z_p(n))\right|
 \leq k(2n-2).
 \tag{SO.7}
\]
This is the direct scheme-theoretic comparison with the
integral Tate cohomology controlled by the Selmer certificates.
The mechanism of Soul\'e's proof is torsion homology of
arithmetic groups, a Hermitian reduction complex, and the
Hurewicz and Chern-class comparisons. It supplies no
vanishing theorem for the groups in (SO.7).
